\chapter*{Bibliography}
\addcontentsline{toc}{chapter}{Bibliography}
\markboth{Bibliography}{Bibliography}
\small
\begin{list}{}{\leftmargin=1.5em\itemindent=-1.5em\itemsep=0.4ex}
\item Abadie, A. and Gardeazabal, J. (2003). The economic costs of conflict: a case study of the Basque Country. \emph{American Economic Review} 93.
\item Abadie, A., Diamond, A. and Hainmueller, J. (2010). Synthetic control methods for comparative case studies. \emph{Journal of the American Statistical Association} 105.
\item Abdulkadiro\u{g}lu, A. and S\"onmez, T. (2003). School choice: a mechanism design approach. \emph{American Economic Review} 93.
\item Acemoglu, D., Johnson, S. and Robinson, J. A. (2001). The colonial origins of comparative development: an empirical investigation. \emph{American Economic Review} 91.
\item Albouy, D. Y. (2012). The colonial origins of comparative development: an empirical investigation: comment. \emph{American Economic Review} 102.
\item Andersen, P. K., Geskus, R. B., de Witte, T. and Putter, H. (2012). Competing risks in epidemiology: possibilities and pitfalls. \emph{International Journal of Epidemiology} 41.
\item Anderson, T. W. and Rubin, H. (1949). Estimation of the parameters of a single equation in a complete system of stochastic equations. \emph{Annals of Mathematical Statistics} 20.
\item Andrews, D. W. K. (1999). Estimation when a parameter is on a boundary. \emph{Econometrica} 67.
\item Angrist, J. D. (1990). Lifetime earnings and the Vietnam era draft lottery: evidence from Social Security administrative records. \emph{American Economic Review} 80.
\item Angrist, J. D. (1998). Estimating the labor market impact of voluntary military service using Social Security data on military applicants. \emph{Econometrica} 66.
\item Angrist, J. D. and Evans, W. N. (1998). Children and their parents' labor supply: evidence from exogenous variation in family size. \emph{American Economic Review} 88.
\item Angrist, J. D. and Imbens, G. W. (1995). Two stage least squares estimation of average causal effects in models with variable treatment intensity. \emph{Journal of the American Statistical Association} 90.
\item Angrist, J. D. and Krueger, A. B. (1991). Does compulsory school attendance affect schooling and earnings? \emph{Quarterly Journal of Economics} 106.
\item Angrist, J. D. and Pischke, J. S. (2010). The credibility revolution in empirical economics: how better research design is taking the con out of econometrics. \emph{Journal of Economic Perspectives} 24.
\item Angrist, J. D., Graddy, K. and Imbens, G. W. (2000). The interpretation of instrumental variables estimators in simultaneous equations models with an application to the demand for fish. \emph{Review of Economic Studies} 67.
\item Angrist, J. D., Imbens, G. W. and Rubin, D. B. (1996). Identification of causal effects using instrumental variables. \emph{Journal of the American Statistical Association} 91.
\item Aranda Ordaz, F. J. (1981). On two families of transformations to additivity for binary response data. \emph{Biometrika} 68.
\item Arkhangelsky, D., Athey, S., Hirshberg, D. A., Imbens, G. W. and Wager, S. (2021). Synthetic difference in differences. \emph{American Economic Review} 111.
\item Armstrong, T. B. (2016). Large market asymptotics for differentiated product demand estimators with economic models of supply. \emph{Econometrica} 84.
\item Aronow, P. M. and Samii, C. (2016). Does regression produce representative estimates of causal effects? \emph{American Journal of Political Science} 60.
\item Aronow, P. M., Green, D. P. and Lee, D. K. K. (2014). Sharp bounds on the variance in randomized experiments. \emph{Annals of Statistics} 42.
\item Ashenfelter, O. and Hosken, D. (2010). The effect of mergers on consumer prices: evidence from five mergers on the enforcement margin. \emph{Journal of Law and Economics} 53.
\item Athey, S. and Imbens, G. W. (2006). Identification and inference in nonlinear difference in differences models. \emph{Econometrica} 74.
\item Athey, S. and Imbens, G. W. (2019). Machine learning methods that economists should know about. \emph{Annual Review of Economics} 11.
\item Athey, S. and Wager, S. (2021). Policy learning with observational data. \emph{Econometrica} 89.
\item Athey, S., Tibshirani, J. and Wager, S. (2019). Generalized random forests. \emph{Annals of Statistics} 47.
\item Audibert, J.-Y. and Tsybakov, A. B. (2007). Fast learning rates for plug-in classifiers. \emph{Annals of Statistics} 35.
\item Aumann, R. J. (1974). Subjectivity and correlation in randomized strategies. \emph{Journal of Mathematical Economics} 1.
\item Aumann, R. J. (1987). Correlated equilibrium as an expression of Bayesian rationality. \emph{Econometrica} 55.
\item Backus, M., Conlon, C. and Sinkinson, M. (2021). Common ownership and competition in the ready to eat cereal industry. NBER Working Paper 28350.
\item Bahadur, R. R. (1967). Rates of convergence of estimates and test statistics. \emph{Annals of Mathematical Statistics} 38.
\item Balakrishnan, S., Kennedy, E. H. and Wasserman, L. (2023). The fundamental limits of structure agnostic functional estimation. arXiv:2305.04116.
\item Balke, A. and Pearl, J. (1997). Bounds on treatment effects from studies with imperfect compliance. \emph{Journal of the American Statistical Association} 92.
\item Barndorff Nielsen, O. E. (1983). On a formula for the distribution of the maximum likelihood estimator. \emph{Biometrika} 70.
\item Basmann, R. L. (1957). A generalized classical method of linear estimation of coefficients in a structural equation. \emph{Econometrica} 25.
\item Basu, D. (1971). An essay on the logical foundations of survey sampling, part one. In \emph{Foundations of Statistical Inference}, Holt, Rinehart and Winston.
\item Battey, H. S. and Edgar, C. (2026). Treatment effect: a critique. arXiv:2601.15467.
\item Beaumont, M. A., Zhang, W. and Balding, D. J. (2002). Approximate Bayesian computation in population genetics. \emph{Genetics} 162.
\item Bellemare, M. G., Dabney, W. and Munos, R. (2017). A distributional perspective on reinforcement learning. \emph{Proceedings of the 34th International Conference on Machine Learning}.
\item Belloni, A., Chen, D., Chernozhukov, V. and Hansen, C. (2012). Sparse models and methods for optimal instruments with an application to eminent domain. \emph{Econometrica} 80.
\item Belloni, A., Chernozhukov, V. and Hansen, C. (2014). Inference on treatment effects after selection among high dimensional controls. \emph{Review of Economic Studies} 81.
\item Bernardo, J. M. (1979). Reference posterior distributions for Bayesian inference. \emph{Journal of the Royal Statistical Society, Series B} 41.
\item Berry, S. (1994). Estimating discrete choice models of product differentiation. \emph{RAND Journal of Economics} 25.
\item Berry, S. and Haile, P. (2014). Identification in differentiated products markets using market level data. \emph{Econometrica} 82.
\item Berry, S., Levinsohn, J. and Pakes, A. (1995). Automobile prices in market equilibrium. \emph{Econometrica} 63.
\item Bertrand, M., Duflo, E. and Mullainathan, S. (2004). How much should we trust differences in differences estimates? \emph{Quarterly Journal of Economics} 119.
\item Bhadra, A., Datta, J., Polson, N. G. and Willard, B. (2017). The horseshoe+ estimator of ultra sparse signals. \emph{Bayesian Analysis} 12.
\item Bhadra, A., Datta, J., Polson, N. G. and Willard, B. (2019). Lasso meets horseshoe: a survey. \emph{Statistical Science} 34.
\item Bhattacharya, A., Chakraborty, A. and Mallick, B. K. (2016). Fast sampling with Gaussian scale mixture priors in high dimensional regression. \emph{Biometrika} 103.
\item Bickel, P. J., Klaassen, C. A. J., Ritov, Y. and Wellner, J. A. (1993). \emph{Efficient and Adaptive Estimation for Semiparametric Models}. Johns Hopkins University Press.
\item Bickel, P. J., Ritov, Y. and Tsybakov, A. B. (2009). Simultaneous analysis of Lasso and Dantzig selector. \emph{Annals of Statistics} 37.
\item Billingsley, P. (1995). \emph{Probability and Measure}, third edition. Wiley.
\item Bingham, N. H., Goldie, C. M. and Teugels, J. L. (1987). \emph{Regular Variation}. Cambridge University Press.
\item Borusyak, K., Jaravel, X. and Spiess, J. (2024). Revisiting event study designs: robust and efficient estimation. \emph{Review of Economic Studies} 91.
\item Bound, J., Jaeger, D. A. and Baker, R. M. (1995). Problems with instrumental variables estimation when the correlation between the instruments and the endogenous explanatory variable is weak. \emph{Journal of the American Statistical Association} 90.
\item Bowden, J., Davey Smith, G. and Burgess, S. (2015). Mendelian randomization with invalid instruments: effect estimation and bias detection through Egger regression. \emph{International Journal of Epidemiology} 44.
\item Bowden, J., Davey Smith, G., Haycock, P. C. and Burgess, S. (2016). Consistent estimation in Mendelian randomization with some invalid instruments using a weighted median estimator. \emph{Genetic Epidemiology} 40.
\item Branson, Z., Rischard, M., Bornn, L. and Miratrix, L. W. (2019). A nonparametric Bayesian methodology for regression discontinuity designs. \emph{Journal of Statistical Planning and Inference} 202.
\item Brenier, Y. (1991). Polar factorization and monotone rearrangement of vector valued functions. \emph{Communications on Pure and Applied Mathematics} 44.
\item Bresnahan, T. (1982). The oligopoly solution concept is identified. \emph{Economics Letters} 10.
\item Breunig, C., Liu, R. and Yu, Z. (2022). Double robust Bayesian inference on average treatment effects. arXiv:2211.16298.
\item Brillinger, D. R. (2012). A generalized linear model with Gaussian regressor variables. In \emph{Selected Works of David Brillinger}, Springer.
\item Brodersen, K. H., Gallusser, F., Koehler, J., Remy, N. and Scott, S. L. (2015). Inferring causal impact using Bayesian structural time series models. \emph{Annals of Applied Statistics} 9.
\item Brown, L. D., George, E. I. and Xu, X. (2008). Admissible predictive density estimation. \emph{Annals of Statistics} 36.
\item Buckles, K. S. and Hungerman, D. M. (2013). Season of birth and later outcomes: old questions, new answers. \emph{Review of Economics and Statistics} 95.
\item Cai, T. T. and Guo, Z. (2017). Confidence intervals for high dimensional linear regression: minimax rates and adaptivity. \emph{Annals of Statistics} 45.
\item Callaway, B. and Sant'Anna, P. H. C. (2021). Difference in differences with multiple time periods. \emph{Journal of Econometrics} 225.
\item Caponnetto, A. and De Vito, E. (2007). Optimal rates for the regularized least squares algorithm. \emph{Foundations of Computational Mathematics} 7.
\item Card, D. (1995). Using geographic variation in college proximity to estimate the return to schooling. In \emph{Aspects of Labour Market Behaviour: Essays in Honour of John Vanderkamp}, University of Toronto Press.
\item Carvalho, C. M., Polson, N. G. and Scott, J. G. (2010). The horseshoe estimator for sparse signals. \emph{Biometrika} 97.
\item Castillo, I. (2012). A semiparametric Bernstein von Mises theorem for Gaussian process priors. \emph{Probability Theory and Related Fields} 152.
\item Castillo, I. and Rousseau, J. (2015). A Bernstein von Mises theorem for smooth functionals in semiparametric models. \emph{Annals of Statistics} 43.
\item Castillo, I., Schmidt Hieber, J. and van der Vaart, A. W. (2015). Bayesian linear regression with sparse priors. \emph{Annals of Statistics} 43.
\item Chamberlain, G. and Imbens, G. W. (2004). Random effects estimators with many instrumental variables. \emph{Econometrica} 72.
\item Charig, C. R., Webb, D. R., Payne, S. R. and Wickham, J. E. A. (1986). Comparison of treatment of renal calculi by open surgery, percutaneous nephrolithotomy, and extracorporeal shockwave lithotripsy. \emph{British Medical Journal} 292.
\item Cheng, C. and Hoekstra, M. (2013). Does strengthening self defense law deter crime or escalate violence? Evidence from expansions to castle doctrine. \emph{Journal of Human Resources} 48.
\item Chernozhukov, V. and Hansen, C. (2004). The effects of 401(k) participation on the wealth distribution: an instrumental quantile regression analysis. \emph{Review of Economics and Statistics} 86.
\item Chernozhukov, V. and Hong, H. (2003). An MCMC approach to classical estimation. \emph{Journal of Econometrics} 115.
\item Chernozhukov, V., Chetverikov, D., Demirer, M., Duflo, E., Hansen, C., Newey, W. K. and Robins, J. (2018). Double/debiased machine learning for treatment and structural parameters. \emph{Econometrics Journal} 21.
\item Chernozhukov, V., Cinelli, C., Newey, W. K., Sharma, A. and Syrgkanis, V. (2022a). Long story short: omitted variable bias in causal machine learning. NBER Working Paper 30302.
\item Chernozhukov, V., Escanciano, J. C., Ichimura, H., Newey, W. K. and Robins, J. M. (2022b). Locally robust semiparametric estimation. \emph{Econometrica} 90.
\item Chernozhukov, V., Fern\'andez Val, I. and Galichon, A. (2010). Quantile and probability curves without crossing. \emph{Econometrica} 78.
\item Chernozhukov, V., Fern\'andez Val, I. and Melly, B. (2013). Inference on counterfactual distributions. \emph{Econometrica} 81.
\item Chernozhukov, V., Galichon, A., Hallin, M. and Henry, M. (2017). Monge Kantorovich depth, quantiles, ranks and signs. \emph{Annals of Statistics} 45.
\item Chiang, H. D., Kato, K., Ma, Y. and Sasaki, Y. (2022). Multiway cluster robust double/debiased machine learning. \emph{Journal of Business and Economic Statistics} 40.
\item Chib, S. and Hamilton, B. H. (2002). Semiparametric Bayes analysis of longitudinal data treatment models. \emph{Journal of Econometrics} 110.
\item Chickering, D. M. (2002). Optimal structure identification with greedy search. \emph{Journal of Machine Learning Research} 3.
\item Chickering, D. M. and Pearl, J. (1996). A clinician's tool for analyzing non compliance. \emph{Proceedings of the Thirteenth National Conference on Artificial Intelligence}, AAAI Press.
\item Chipman, H. A., George, E. I. and McCulloch, R. E. (2010). BART: Bayesian additive regression trees. \emph{Annals of Applied Statistics} 4.
\item Cinelli, C. and Hazlett, C. (2020). Making sense of sensitivity: extending omitted variable bias. \emph{Journal of the Royal Statistical Society, Series B} 82.
\item Clarke, P. S. and Polselli, A. (2025). Double machine learning for static panel models with fixed effects. \emph{Econometrics Journal}.
\item Claxton, K. (1999). The irrelevance of inference: a decision making approach to the stochastic evaluation of health care technologies. \emph{Journal of Health Economics} 18.
\item Conley, T. G., Hansen, C. B. and Rossi, P. E. (2012). Plausibly exogenous. \emph{Review of Economics and Statistics} 94.
\item Conley, T. G., Hansen, C. B., McCulloch, R. E. and Rossi, P. E. (2008). A semiparametric Bayesian approach to the instrumental variable problem. \emph{Journal of Econometrics} 144.
\item Conlon, C. and Gortmaker, J. (2020). Best practices for differentiated products demand estimation with PyBLP. \emph{RAND Journal of Economics} 51.
\item Cook, S. R., Gelman, A. and Rubin, D. B. (2006). Validation of software for Bayesian models using posterior quantiles. \emph{Journal of Computational and Graphical Statistics} 15.
\item Cornfield, J., Haenszel, W., Hammond, E. C., Lilienfeld, A. M., Shimkin, M. B. and Wynder, E. L. (1959). Smoking and lung cancer: recent evidence and a discussion of some questions. \emph{Journal of the National Cancer Institute} 22.
\item Cox, D. R. (1958). \emph{Planning of Experiments}. Wiley.
\item Cox, D. R. (1992). Causality: some statistical aspects. \emph{Journal of the Royal Statistical Society, Series A} 155.
\item Cox, D. R. and Wermuth, N. (1996). \emph{Multivariate Dependencies: Models, Analysis and Interpretation}. Chapman and Hall.
\item Cox, D. R. and Wermuth, N. (2004). Causality: a statistical view. \emph{International Statistical Review} 72.
\item Crump, R. K., Hotz, V. J., Imbens, G. W. and Mitnik, O. A. (2009). Dealing with limited overlap in estimation of average treatment effects. \emph{Biometrika} 96.
\item Csisz\'ar, I. (1975). I divergence geometry of probability distributions and minimization problems. \emph{Annals of Probability} 3.
\item Csisz\'ar, I. (1984). Sanov property, generalized I projection and a conditional limit theorem. \emph{Annals of Probability} 12.
\item Cunningham, S. (2021). \emph{Causal Inference: The Mixtape}. Yale University Press.
\item Dabney, W., Ostrovski, G., Silver, D. and Munos, R. (2018). Implicit quantile networks for distributional reinforcement learning. \emph{Proceedings of Machine Learning Research} 80.
\item Datta, J. and Polson, N. G. (2023). Quantile importance sampling. arXiv:2305.03158.
\item Datta, J. and Polson, N. G. (2025). Prediction powered inference with inverse probability weighting. arXiv:2508.10149.
\item Datta, J. and Polson, N. G. (2026). Inverse probability weighting: from survey sampling to evidence estimation. \emph{The New England Journal of Statistics in Data Science} 4.
\item \emph{Daubert v. Merrell Dow Pharmaceuticals, Inc.} (1995). 43 F.3d 1311 (9th Cir.).
\item Davey Smith, G. and Ebrahim, S. (2003). Mendelian randomization: can genetic epidemiology contribute to understanding environmental determinants of disease? \emph{International Journal of Epidemiology} 32.
\item Dawid, A. P. (2000). Causal inference without counterfactuals. \emph{Journal of the American Statistical Association} 95.
\item Dawid, A. P. (2002). Influence diagrams for causal modelling and inference. \emph{International Statistical Review} 70.
\item Dawid, A. P. (2010). Beware of the DAG! \emph{Proceedings of the NIPS 2008 Workshop on Causality}, JMLR Workshop and Conference Proceedings 6.
\item Dawid, A. P. (2015). Statistical causality from a decision theoretic perspective. \emph{Annual Review of Statistics and Its Application} 2.
\item Dawid, A. P. (2021). Decision theoretic foundations for statistical causality. \emph{Journal of Causal Inference} 9.
\item de Chaisemartin, C. and D'Haultf\oe uille, X. (2020). Two way fixed effects estimators with heterogeneous treatment effects. \emph{American Economic Review} 110.
\item de Finetti, B. (1937). La pr\'evision: ses lois logiques, ses sources subjectives. \emph{Annales de l'Institut Henri Poincar\'e} 7.
\item Dehejia, R. H. and Wahba, S. (1999). Causal effects in nonexperimental studies: reevaluating the evaluation of training programs. \emph{Journal of the American Statistical Association} 94.
\item Dehejia, R. H. and Wahba, S. (2002). Propensity score-matching methods for nonexperimental causal studies. \emph{Review of Economics and Statistics} 84.
\item Dhaene, J., Kukush, A., Linders, D. and Tang, Q. (2012). Remarks on quantiles and distortion risk measures. \emph{European Actuarial Journal} 2.
\item Diaconis, P. and Freedman, D. (1980). Finite exchangeable sequences. \emph{Annals of Probability} 8.
\item Ding, P. and VanderWeele, T. J. (2016). Sensitivity analysis without assumptions. \emph{Epidemiology} 27.
\item DiTraglia, F. J. and Liu, L. (2025). Bayesian double machine learning for causal inference. arXiv:2508.12688.
\item Doll, R. and Hill, A. B. (1950). Smoking and carcinoma of the lung. \emph{British Medical Journal} 2.
\item Doll, R. and Hill, A. B. (1954). The mortality of doctors in relation to their smoking habits. \emph{British Medical Journal} 1.
\item Donoho, D. L., Johnstone, I. M., Hoch, J. C. and Stern, A. S. (1992). Maximum entropy and the nearly black object. \emph{Journal of the Royal Statistical Society, Series B} 54.
\item Donsker, M. D. and Varadhan, S. R. S. (1975). Asymptotic evaluation of certain Markov process expectations for large time, I. \emph{Communications on Pure and Applied Mathematics} 28.
\item Dorie, V., Hill, J., Shalit, U., Scott, M. and Cervone, D. (2019). Automated versus do it yourself methods for causal inference: lessons learned from a data analysis competition. \emph{Statistical Science} 34.
\item Duarte, M., Magnolfi, L., S{\o}lvsten, M. and Sullivan, C. (2024). Testing firm conduct. \emph{Quantitative Economics} 15.
\item Dub\'e, J. P., Fox, J. T. and Su, C. L. (2012). Improving the numerical performance of static and dynamic aggregate discrete choice random coefficients demand estimation. \emph{Econometrica} 80.
\item Duflo, E., Greenstone, M., Pande, R. and Ryan, N. (2013). Truth-telling by third-party auditors and the response of polluting firms: experimental evidence from India. \emph{Quarterly Journal of Economics} 128.
\item Dufour, J.-M. (1997). Some impossibility theorems in econometrics with applications to structural and dynamic models. \emph{Econometrica} 65.
\item Efron, B. (2011). Tweedie's formula and selection bias. \emph{Journal of the American Statistical Association} 106.
\item Enikolopov, R., Korovkin, V., Petrova, M., Sonin, K. and Zakharov, A. (2013). Field experiment estimate of electoral fraud in Russian parliamentary elections. \emph{Proceedings of the National Academy of Sciences} 110.
\item Enke, B. and Zimmermann, F. (2019). Correlation neglect in belief formation. \emph{Review of Economic Studies} 86.
\item Ericson, W. A. (1969). Subjective Bayesian models in sampling finite populations. \emph{Journal of the Royal Statistical Society, Series B} 31.
\item Farrell, J. and Shapiro, C. (2010). Antitrust evaluation of horizontal mergers: an economic alternative to market definition. \emph{The B.E. Journal of Theoretical Economics} 10.
\item Farrell, M. H., Liang, T. and Misra, S. (2021). Deep neural networks for estimation and inference. \emph{Econometrica} 89.
\item Fearnhead, P. and Prangle, D. (2012). Constructing summary statistics for approximate Bayesian computation: semi automatic approximate Bayesian computation. \emph{Journal of the Royal Statistical Society, Series B} 74.
\item Ferguson, T. S. (1973). A Bayesian analysis of some nonparametric problems. \emph{Annals of Statistics} 1.
\item Firpo, S. (2007). Efficient semiparametric estimation of quantile treatment effects. \emph{Econometrica} 75.
\item Fisher, R. A. (1925). \emph{Statistical Methods for Research Workers}. Oliver and Boyd.
\item Fisher, R. A. (1934). Two new properties of mathematical likelihood. \emph{Proceedings of the Royal Society of London, Series A} 144.
\item Fisher, R. A. (1935). \emph{The Design of Experiments}. Oliver and Boyd.
\item Fisher, R. A. (1957). Dangers of cigarette smoking. \emph{British Medical Journal} 2.
\item Fisher, R. A. (1958). Cancer and smoking. \emph{Nature} 182.
\item Fisher, R. A. (1959). \emph{Smoking: The Cancer Controversy}. Oliver and Boyd.
\item Fisher, R. A. and Bartlett, S. (1931). Pasteurized and raw milk. \emph{Nature} 127.
\item Fort, M. K. (1955). The embedding of homeomorphisms in flows. \emph{Proceedings of the American Mathematical Society} 6.
\item Frandsen, B., Lefgren, L. and Leslie, E. (2023). Judging judge fixed effects. \emph{American Economic Review} 113.
\item Frangakis, C. E. and Rubin, D. B. (2002). Principal stratification in causal inference. \emph{Biometrics} 58.
\item Fraser, D. A. S. (1961). The fiducial method and invariance. \emph{Biometrika} 48.
\item Frazier, D. T., Martin, G. M., Robert, C. P. and Rousseau, J. (2018). Asymptotic properties of approximate Bayesian computation. \emph{Biometrika} 105.
\item Frisch, R. (1934). \emph{Statistical Confluence Analysis by Means of Complete Regression Systems}. University Institute of Economics, Oslo.
\item Frisch, R. and Waugh, F. V. (1933). Partial time regressions as compared with individual trends. \emph{Econometrica} 1.
\item Gail, M. H., Wieand, S. and Piantadosi, S. (1984). Biased estimates of treatment effect in randomized experiments with nonlinear regressions and omitted covariates. \emph{Biometrika} 71.
\item Gandhi, A. and Houde, J. F. (2019). Measuring substitution patterns in differentiated products industries. NBER Working Paper 26375.
\item Ganesh, A. J. and O'Connell, N. (1999). An inverse of Sanov's theorem. \emph{Statistics and Probability Letters} 42.
\item Geary, R. C. (1949). Determination of linear relations between systematic parts of variables with errors of observation the variances of which are unknown. \emph{Econometrica} 17.
\item Gelman, A., Carlin, J. B., Stern, H. S., Dunson, D. B., Vehtari, A. and Rubin, D. B. (2013). \emph{Bayesian Data Analysis}, third edition. Chapman and Hall/CRC.
\item Gelman, A. and Hill, J. (2007). \emph{Data Analysis Using Regression and Multilevel/Hierarchical Models}. Cambridge University Press.
\item Gelman, A., Meng, X.-L. and Stern, H. S. (1996). Posterior predictive assessment of model fitness via realized discrepancies. \emph{Statistica Sinica} 6.
\item Gelman, A., Vehtari, A., Simpson, D., Margossian, C. C., Carpenter, B., Yao, Y., Kennedy, L., Gabry, J., B\"urkner, P.-C. and Modr\'ak, M. (2020). Bayesian workflow. arXiv:2011.01808.
\item George, E. I., Liang, F. and Xu, X. (2006). Improved minimax predictive densities under Kullback Leibler loss. \emph{Annals of Statistics} 34.
\item Ghosal, S. and van der Vaart, A. W. (2017). \emph{Fundamentals of Nonparametric Bayesian Inference}. Cambridge University Press.
\item Ghosal, S., Ghosh, J. K. and van der Vaart, A. W. (2000). Convergence rates of posterior distributions. \emph{Annals of Statistics} 28.
\item Ghosh, J. K. (2015). Weak paradoxes and paradigms. In \emph{Statistical Paradigms: Recent Advances and Reconciliations}, World Scientific.
\item Gibbard, A. and Harper, W. L. (1978). Counterfactuals and two kinds of expected utility. In \emph{Foundations and Applications of Decision Theory}, Reidel.
\item Gin\'e, E., G\"otze, F. and Mason, D. M. (1997). When is the Student $t$ statistic asymptotically standard normal? \emph{Annals of Probability} 25.
\item Gkatzionis, A. and Burgess, S. (2019). Contextualizing selection bias in Mendelian randomization: how bad is it likely to be? \emph{International Journal of Epidemiology} 48.
\item Gleser, L. J. and Hwang, J. T. (1987). The nonexistence of $100(1-\alpha)\%$ confidence sets of finite expected diameter in errors in variables and related models. \emph{Annals of Statistics} 15.
\item Gneiting, T. and Raftery, A. E. (2007). Strictly proper scoring rules, prediction, and estimation. \emph{Journal of the American Statistical Association} 102.
\item Goodman Bacon, A. (2021). Difference in differences with variation in treatment timing. \emph{Journal of Econometrics} 225.
\item Graddy, K. (1995). Testing for imperfect competition at the Fulton fish market. \emph{RAND Journal of Economics} 26.
\item Greenland, S., Robins, J. M. and Pearl, J. (1999). Confounding and collapsibility in causal inference. \emph{Statistical Science} 14.
\item Gron, A., Polson, N. G. and Viard, V. B. (2003). Likelihood based estimation of nonlinear equilibrium models with random coefficients. Working paper, Kellogg School of Management, Northwestern University.
\item Guo, Z., Wang, W., Cai, T. T. and Li, H. (2019). Optimal estimation of genetic relatedness in high dimensional linear models. \emph{Journal of the American Statistical Association} 114.
\item Gustafson, P. (2015). \emph{Bayesian Inference for Partially Identified Models}. Chapman and Hall/CRC.
\item Haavelmo, T. (1943). The statistical implications of a system of simultaneous equations. \emph{Econometrica} 11.
\item Haavelmo, T. (1944). The probability approach in econometrics. \emph{Econometrica} 12, supplement.
\item Hahn, J. (1998). On the role of the propensity score in efficient semiparametric estimation of average treatment effects. \emph{Econometrica} 66.
\item Hahn, P. R., Carvalho, C. M., Puelz, D. and He, J. (2018). Regularization and confounding in linear regression for treatment effect estimation. \emph{Bayesian Analysis} 13.
\item Hahn, P. R., Murray, J. S. and Carvalho, C. M. (2020). Bayesian regression tree models for causal inference. \emph{Bayesian Analysis} 15.
\item Hainmueller, J. (2012). Entropy balancing for causal effects. \emph{Political Analysis} 20.
\item H\'ajek, J. (1971). Comment on ``An essay on the logical foundations of survey sampling'' by D. Basu. In \emph{Foundations of Statistical Inference}, Holt, Rinehart and Winston.
\item Hammond, E. C. (1964). Smoking in relation to mortality and morbidity: findings in first thirty four months of follow up in a prospective study started in 1959. \emph{Journal of the National Cancer Institute} 32.
\item Hammond, E. C. and Horn, D. (1958). Smoking and death rates: report on forty four months of follow up of 187,783 men. II. Death rates by cause. \emph{Journal of the American Medical Association} 166.
\item Hammond, L., Fox, J., Everitt, T., Carey, R., Abate, A. and Wooldridge, M. (2023). Reasoning about causality in games. \emph{Artificial Intelligence} 320.
\item Hardy, G. H., Littlewood, J. E. and P\'olya, G. (1952). \emph{Inequalities}, second edition. Cambridge University Press.
\item Harmeling, S. and Toussaint, M. (2007). Bayesian estimators for Robins and Ritov's problem. Technical report.
\item Harrell, F. E. (2026). Causal by design. Statistical Thinking blog, fharrell.com/post/causal.
\item Hausman, J., Leonard, G. and Zona, J. D. (1994). Competitive analysis with differentiated products. \emph{Annales d'\'Economie et de Statistique} 34.
\item He, J. and Hahn, P. R. (2023). Stochastic tree ensembles for regularized nonlinear regression. \emph{Journal of the American Statistical Association} 118.
\item He, J., Yalov, S. and Hahn, P. R. (2019). XBART: accelerated Bayesian additive regression trees. \emph{Proceedings of the 22nd International Conference on Artificial Intelligence and Statistics}, PMLR 89.
\item Heath, D. and Sudderth, W. (1976). De Finetti's theorem on exchangeable variables. \emph{American Statistician} 30.
\item Heckman, J. J. (1979). Sample selection bias as a specification error. \emph{Econometrica} 47.
\item Heckman, J. J. and Vytlacil, E. J. (2005). Structural equations, treatment effects, and econometric policy evaluation. \emph{Econometrica} 73.
\item Heckman, J. J. and Vytlacil, E. J. (2007). Econometric evaluation of social programs, part I. In \emph{Handbook of Econometrics} 6B, Elsevier.
\item Heckman, J. J., Smith, J. A. and Clements, N. (1997). Making the most out of programme evaluations and social experiments: accounting for heterogeneity in programme impacts. \emph{Review of Economic Studies} 64.
\item Hern\'an, M. A. and Robins, J. M. (2016). Using big data to emulate a target trial when a randomized trial is not available. \emph{American Journal of Epidemiology} 183.
\item Hern\'an, M. A. and Robins, J. M. (2020). \emph{Causal Inference: What If}. Chapman and Hall/CRC.
\item Hewitt, E. and Savage, L. J. (1955). Symmetric measures on Cartesian products. \emph{Transactions of the American Mathematical Society} 80.
\item Hill, A. B. (1965). The environment and disease: association or causation? \emph{Proceedings of the Royal Society of Medicine} 58.
\item Hill, J. L. (2011). Bayesian nonparametric modeling for causal inference. \emph{Journal of Computational and Graphical Statistics} 20.
\item Hill, J. L. and Su, Y. S. (2013). Assessing lack of common support in causal inference using Bayesian nonparametrics. \emph{Annals of Applied Statistics} 7.
\item Hill, J. L., Linero, A. R. and Murray, J. S. (2020). Bayesian additive regression trees: a review and look forward. \emph{Annual Review of Statistics and Its Application} 7.
\item Hirano, K., Imbens, G. W. and Ridder, G. (2003). Efficient estimation of average treatment effects using the estimated propensity score. \emph{Econometrica} 71.
\item Hjelmborg, J. et al. (2017). Lung cancer, genetic predisposition and smoking: the Nordic Twin Study of Cancer. \emph{Thorax} 72.
\item Hoeffding, W. (1965). Asymptotically optimal tests for multinomial distributions. \emph{Annals of Mathematical Statistics} 36.
\item Hoerl, A. E. and Kennard, R. W. (1970). Ridge regression: biased estimation for nonorthogonal problems. \emph{Technometrics} 12.
\item Hoffman, M. D. and Gelman, A. (2014). The No-U-Turn Sampler: adaptively setting path lengths in Hamiltonian Monte Carlo. \emph{Journal of Machine Learning Research} 15.
\item Holland, P. W. (1986). Statistics and causal inference. \emph{Journal of the American Statistical Association} 81.
\item Hora, R. B. and Buehler, R. J. (1966). Fiducial theory and invariant estimation. \emph{Annals of Mathematical Statistics} 37.
\item Horvitz, D. G. and Thompson, D. J. (1952). A generalization of sampling without replacement from a finite universe. \emph{Journal of the American Statistical Association} 47.
\item Hoxby, C. M. (2000). Does competition among public schools benefit students and taxpayers? \emph{American Economic Review} 90.
\item Huang, Y. and Valtorta, M. (2006). Pearl's calculus of intervention is complete. \emph{Proceedings of the 22nd Conference on Uncertainty in Artificial Intelligence}.
\item Hudgens, M. G. and Halloran, M. E. (2008). Toward causal inference with interference. \emph{Journal of the American Statistical Association} 103.
\item Hung, R. J. et al. (2008). A susceptibility locus for lung cancer maps to nicotinic acetylcholine receptor subunit genes on 15q25. \emph{Nature} 452.
\item Huntington Klein, N. (2021). \emph{The Effect: An Introduction to Research Design and Causality}. Chapman and Hall/CRC.
\item ICH (2019). E9(R1) addendum on estimands and sensitivity analysis in clinical trials to the guideline on statistical principles for clinical trials. International Council for Harmonisation.
\item Imai, K. and Ratkovic, M. (2014). Covariate balancing propensity score. \emph{Journal of the Royal Statistical Society, Series B} 76.
\item Imbens, G. W. (2020). Potential outcome and directed acyclic graph approaches to causality: relevance for empirical practice in economics. \emph{Journal of Economic Literature} 58.
\item Imbens, G. W. and Angrist, J. D. (1994). Identification and estimation of local average treatment effects. \emph{Econometrica} 62.
\item Imbens, G. W. and Lemieux, T. (2008). Regression discontinuity designs: a guide to practice. \emph{Journal of Econometrics} 142.
\item Imbens, G. W. and Rubin, D. B. (1997). Bayesian inference for causal effects in randomized experiments with noncompliance. \emph{Annals of Statistics} 25.
\item Imbens, G. W. and Rubin, D. B. (2015). \emph{Causal Inference for Statistics, Social, and Biomedical Sciences}. Cambridge University Press.
\item James, W. and Stein, C. (1961). Estimation with quadratic loss. \emph{Proceedings of the Fourth Berkeley Symposium on Mathematical Statistics and Probability} 1.
\item Javanmard, A. and Montanari, A. (2014). Confidence intervals and hypothesis testing for high dimensional regression. \emph{Journal of Machine Learning Research} 15.
\item Jiang, B., Wu, T.-Y., Zheng, C. and Wong, W. H. (2017). Learning summary statistic for approximate Bayesian computation via deep neural network. \emph{Statistica Sinica} 27.
\item Jiang, R., Manchanda, P. and Rossi, P. E. (2009). Bayesian analysis of random coefficient logit models using aggregate data. \emph{Journal of Econometrics} 149.
\item Joyce, J. M. (1999). \emph{The Foundations of Causal Decision Theory}. Cambridge University Press.
\item Julious, S. A. and Mullee, M. A. (1994). Confounding and Simpson's paradox. \emph{British Medical Journal} 309.
\item Kadane, J. B. and Seidenfeld, T. (1990). Randomization in a Bayesian perspective. \emph{Journal of Statistical Planning and Inference} 25.
\item Kahneman, D. and Tversky, A. (1979). Prospect theory: an analysis of decision under risk. \emph{Econometrica} 47.
\item Kallenberg, O. (2002). \emph{Foundations of Modern Probability}, second edition. Springer.
\item Kallus, N. (2022). What's the harm? Sharp bounds on the fraction negatively affected by treatment. \emph{Advances in Neural Information Processing Systems} 35.
\item Kallus, N., Mao, X. and Uehara, M. (2024). Localized debiased machine learning: efficient inference on quantile treatment effects and beyond. \emph{Journal of Machine Learning Research} 25.
\item Kang, J. D. Y. and Schafer, J. L. (2007). Demystifying double robustness. \emph{Statistical Science} 22.
\item Katan, M. B. (1986). Apolipoprotein E isoforms, serum cholesterol, and cancer. \emph{The Lancet} 327.
\item Kennedy, E. H. (2023). Towards optimal doubly robust estimation of heterogeneous causal effects. \emph{Electronic Journal of Statistics} 17.
\item Khan, S. and Ugander, J. (2023). Adaptive normalization for IPW estimation. \emph{Journal of Causal Inference} 11.
\item Kleibergen, F. and Zivot, E. (2003). Bayesian and classical approaches to instrumental variable regression. \emph{Journal of Econometrics} 114.
\item Kleijn, B. J. K. and van der Vaart, A. W. (2012). The Bernstein von Mises theorem under misspecification. \emph{Electronic Journal of Statistics} 6.
\item Kling, J. R. (2006). Incarceration length, employment, and earnings. \emph{American Economic Review} 96.
\item Knapik, B. T., Szab\'o, B. T., van der Vaart, A. W. and van Zanten, J. H. (2016). Bayes procedures for adaptive inference in inverse problems for the white noise model. \emph{Probability Theory and Related Fields} 164.
\item Knapik, B. T., van der Vaart, A. W. and van Zanten, J. H. (2011). Bayesian inverse problems with Gaussian priors. \emph{Annals of Statistics} 39.
\item Knittel, C. R. and Metaxoglou, K. (2014). Estimation of random coefficient demand models: two empiricists' perspective. \emph{Review of Economics and Statistics} 96.
\item Knothe, H. (1957). Contributions to the theory of convex bodies. \emph{Michigan Mathematical Journal} 4.
\item Koenker, R. and Bassett, G. (1978). Regression quantiles. \emph{Econometrica} 46.
\item Koller, D. and Milch, B. (2003). Multi agent influence diagrams for representing and solving games. \emph{Games and Economic Behavior} 45.
\item Koopmans, T. C., Rubin, H. and Leipnik, R. B. (1950). Measuring the equation systems of dynamic economics. In \emph{Statistical Inference in Dynamic Economic Models}, Cowles Commission Monograph 10, Wiley.
\item Kopell, N. (1970). Commuting diffeomorphisms. In \emph{Global Analysis}, Proceedings of Symposia in Pure Mathematics 14, American Mathematical Society.
\item Krinsky, I. and Robb, A. L. (1986). On approximating the statistical properties of elasticities. \emph{Review of Economics and Statistics} 68.
\item K\"unzel, S. R., Sekhon, J. S., Bickel, P. J. and Yu, B. (2019). Metalearners for estimating heterogeneous treatment effects using machine learning. \emph{Proceedings of the National Academy of Sciences} 116.
\item Kupper, M. and Schachermayer, W. (2009). Representation results for law invariant time consistent functions. \emph{Mathematics and Financial Economics} 2.
\item Laibson, D. (1997). Golden eggs and hyperbolic discounting. \emph{Quarterly Journal of Economics} 112.
\item Lakshminarayanan, B., Roy, D. M. and Teh, Y. W. (2015). Particle Gibbs for Bayesian additive regression trees. \emph{Proceedings of the 18th International Conference on Artificial Intelligence and Statistics}, PMLR 38.
\item LaLonde, R. J. (1986). Evaluating the econometric evaluations of training programs with experimental data. \emph{American Economic Review} 76.
\item Larsson, S. C., Carter, P., Kar, S., Vithayathil, M., Mason, A. M., Micha\"elsson, K. and Burgess, S. (2020). Smoking, alcohol consumption, and cancer: a mendelian randomisation study in UK Biobank and international genetic consortia participants. \emph{PLOS Medicine} 17.
\item Lee, D. S. (2008). Randomized experiments from non random selection in U.S. House elections. \emph{Journal of Econometrics} 142.
\item Lei, L. and Cand\`es, E. J. (2021). Conformal inference of counterfactuals and individual treatment effects. \emph{Journal of the Royal Statistical Society, Series B} 83.
\item Leighton, G. and McKinlay, P. L. (1930). \emph{Milk Consumption and the Growth of School Children}. Department of Health for Scotland, HMSO.
\item Levitt, S. D. (1997). Using electoral cycles in police hiring to estimate the effect of police on crime. \emph{American Economic Review} 87.
\item Lewis, D. (1973). Causation. \emph{Journal of Philosophy} 70.
\item Lewis, D. (1981). Causal decision theory. \emph{Australasian Journal of Philosophy} 59.
\item Li, F., Morgan, K. L. and Zaslavsky, A. M. (2018). Balancing covariates via propensity score weighting. \emph{Journal of the American Statistical Association} 113.
\item Li, L. (2010). Are Bayesian inferences weak for Wasserman's example? \emph{Communications in Statistics, Simulation and Computation} 39.
\item Lindley, D. V. (1958). Fiducial distributions and Bayes' theorem. \emph{Journal of the Royal Statistical Society, Series B} 20.
\item Lindley, D. V. and Novick, M. R. (1981). The role of exchangeability in inference. \emph{Annals of Statistics} 9.
\item Linero, A. R. (2018). Bayesian regression trees for high dimensional prediction and variable selection. \emph{Journal of the American Statistical Association} 113.
\item Linero, A. R. (2024). In nonparametric and high dimensional models, Bayesian ignorability is an informative prior. \emph{Journal of the American Statistical Association} 119.
\item Linero, A. R. and Yang, Y. (2018). Bayesian regression tree ensembles that adapt to smoothness and sparsity. \emph{Journal of the Royal Statistical Society, Series B} 80.
\item Loewenstein, G., O'Donoghue, T. and Rabin, M. (2003). Projection bias in predicting future utility. \emph{Quarterly Journal of Economics} 118.
\item Lovell, M. C. (1963). Seasonal adjustment of economic time series and multiple regression analysis. \emph{Journal of the American Statistical Association} 58.
\item Lucas, R. E. (1976). Econometric policy evaluation: a critique. \emph{Carnegie Rochester Conference Series on Public Policy} 1.
\item Lueckmann, J.-M., Boelts, J., Greenberg, D. S., Gon\c{c}alves, P. J. and Macke, J. H. (2021). Benchmarking simulation based inference. \emph{Proceedings of the 24th International Conference on Artificial Intelligence and Statistics}, PMLR 130.
\item Madansky, A. (1959). The fitting of straight lines when both variables are subject to error. \emph{Journal of the American Statistical Association} 54.
\item Madansky, A. (1964). Instrumental variables in factor analysis. \emph{Psychometrika} 29.
\item Makalic, E. and Schmidt, D. F. (2016). A simple sampler for the horseshoe estimator. \emph{IEEE Signal Processing Letters} 23.
\item Makarov, G. D. (1981). Estimates for the distribution function of a sum of two random variables when the marginal distributions are fixed. \emph{Theory of Probability and its Applications} 26.
\item Manski, C. F. (1990). Nonparametric bounds on treatment effects. \emph{American Economic Review} 80.
\item Manski, C. F. (2003). \emph{Partial Identification of Probability Distributions}. Springer.
\item McCandless, L. C., Gustafson, P. and Levy, A. (2007). Bayesian sensitivity analysis for unmeasured confounding in observational studies. \emph{Statistics in Medicine} 26.
\item McCrary, J. (2002). Using electoral cycles in police hiring to estimate the effect of police on crime: comment. \emph{American Economic Review} 92.
\item McCullagh, P. (2002). What is a statistical model? \emph{Annals of Statistics} 30.
\item McCullagh, P. (2022). \emph{Ten Projects in Applied Statistics}. Springer.
\item Miguel, E., Satyanath, S. and Sergenti, E. (2004). Economic shocks and civil conflict: an instrumental variables approach. \emph{Journal of Political Economy} 112.
\item Miller, N. H. and Weinberg, M. C. (2017). Understanding the price effects of the MillerCoors joint venture. \emph{Econometrica} 85.
\item Moon, H. R. and Schorfheide, F. (2012). Bayesian and frequentist inference in partially identified models. \emph{Econometrica} 80.
\item Morrow, W. R. and Skerlos, S. J. (2011). Fixed point approaches to computing Bertrand Nash equilibrium prices under mixed logit demand. \emph{Operations Research} 59.
\item Mueller, S., Li, A. and Pearl, J. (2022). Causes of effects: learning individual responses from population data. \emph{Proceedings of IJCAI}.
\item Mukherjee, G. and Johnstone, I. M. (2015). Exact minimax estimation of the predictive density in sparse Gaussian models. \emph{Annals of Statistics} 43.
\item Myerson, R. B. (1979). Incentive compatibility and the bargaining problem. \emph{Econometrica} 47.
\item Nareklishvili, M., Polson, N. G. and Sokolov, V. (2023a). Generative causal inference. arXiv:2306.16096.
\item Nareklishvili, M., Polson, N. G. and Sokolov, V. (2023b). Deep partial least squares for instrumental variable regression. \emph{Applied Stochastic Models in Business and Industry}.
\item Nareklishvili, M., Polson, N. G. and Sokolov, V. (2023c). Feature selection for personalized policy analysis. arXiv:2301.00251.
\item Nareklishvili, M., Polson, N. G. and Sokolov, V. (2025). Generative learner for distributional causal effects. arXiv:2503.17894.
\item Nareklishvili, M., Polson, N. G. and Sokolov, V. (2026a). Double machine learning versus hierarchical Bayes. Working paper.
\item Nareklishvili, M., Polson, N. G. and Sokolov, V. (2026b). Optimal transport for causal inference. Working paper.
\item Nelsen, R. B. (2006). \emph{An Introduction to Copulas}, second edition. Springer.
\item Nevo, A. (2000). Mergers with differentiated products: the case of the ready to eat cereal industry. \emph{RAND Journal of Economics} 31.
\item Nevo, A. (2001). Measuring market power in the ready to eat cereal industry. \emph{Econometrica} 69.
\item Newey, W. K. and Powell, J. L. (2003). Instrumental variable estimation of nonparametric models. \emph{Econometrica} 71.
\item Newton, M. A. and Raftery, A. E. (1994). Approximate Bayesian inference with the weighted likelihood bootstrap. \emph{Journal of the Royal Statistical Society, Series B} 56.
\item Newton, M. A., Polson, N. G. and Xu, J. (2021). Weighted Bayesian bootstrap for scalable posterior distributions. \emph{Canadian Journal of Statistics} 49.
\item Neyman, J. (1923, translated 1990). On the application of probability theory to agricultural experiments. \emph{Statistical Science} 5.
\item Neyman, J. (1959). Optimal asymptotic tests of composite statistical hypotheses. In \emph{Probability and Statistics: The Harald Cram\'er Volume}, Wiley.
\item Neyman, J., Iwaszkiewicz, K. and Kolodziejczyk, S. (1935). Statistical problems in agricultural experimentation. \emph{Supplement to the Journal of the Royal Statistical Society} 2.
\item Nie, X. and Wager, S. (2021). Quasi oracle estimation of heterogeneous treatment effects. \emph{Biometrika} 108.
\item Padilla, O. H. M., Tansey, W. and Chen, Y. (2022). Quantile regression with ReLU networks: estimators and minimax rates. \emph{Journal of Machine Learning Research} 23.
\item Pearl, J. (1995). Causal diagrams for empirical research. \emph{Biometrika} 82.
\item Pearl, J. (1999). Probabilities of causation: three counterfactual interpretations and their identification. \emph{Synthese} 121.
\item Pearl, J. (2001). Bayesianism and causality, or, why I am only a half Bayesian. In \emph{Foundations of Bayesianism}, Kluwer.
\item Pearl, J. (2009). \emph{Causality}, second edition. Cambridge University Press.
\item Pearl, J. (2011). Invited commentary: understanding bias amplification. \emph{American Journal of Epidemiology} 174.
\item Pearl, J. and Bareinboim, E. (2014). External validity: from do calculus to transportability across populations. \emph{Statistical Science} 29.
\item Pearl, J. and Mackenzie, D. (2018). \emph{The Book of Why: The New Science of Cause and Effect}. Basic Books.
\item Petrin, A. (2002). Quantifying the benefits of new products: the case of the minivan. \emph{Journal of Political Economy} 110.
\item Philippe, A. and Robert, C. P. (2001). Riemann sums for MCMC estimation and convergence monitoring. \emph{Statistics and Computing} 11.
\item Piironen, J. and Vehtari, A. (2017). Sparsity information and regularization in the horseshoe and other shrinkage priors. \emph{Electronic Journal of Statistics} 11.
\item Pitman, E. J. G. (1939). The estimation of the location and scale parameters of a continuous population of any given form. \emph{Biometrika} 30.
\item Polson, N. G. and Scott, J. G. (2011). Shrink globally, act locally: sparse Bayesian regularization and prediction. In \emph{Bayesian Statistics 9}, Oxford University Press.
\item Polson, N. G. and Scott, J. G. (2014). Vertical likelihood Monte Carlo. arXiv:1409.3601.
\item Polson, N. G. and Sokolov, V. (2023). Generative AI for Bayesian computation. arXiv:2305.14972.
\item Polson, N. G. and Sokolov, V. (2026a). \emph{Generative Bayesian Computation: Quantile Neural Networks for Inference and Surrogates}. Book manuscript.
\item Polson, N. G. and Sokolov, V. (2026b). Orthogonality is not enough: double machine learning and the Bayesian alternative. Working paper.
\item Polson, N. G. and Sokolov, V. (2026c). \emph{Elements of Bayes and Information Theory: A Predictive Theory of Inference}. Book manuscript.
\item Polson, N. G., Ruggeri, F. and Sokolov, V. (2024). Generative Bayesian computation for maximum expected utility. \emph{Entropy} 26.
\item Polson, N. G., Scott, J. G. and Windle, J. (2013). Bayesian inference for logistic models using P\'olya Gamma latent variables. \emph{Journal of the American Statistical Association} 108.
\item Poterba, J. M., Venti, S. F. and Wise, D. A. (1995). Do 401(k) contributions crowd out other personal saving? \emph{Journal of Public Economics} 58.
\item Rambachan, A. and Roth, J. (2023). A more credible approach to parallel trends. \emph{Review of Economic Studies} 90.
\item Rathouz, P. J. (2004). Fixed effect models for longitudinal binary data with drop outs missing at random. \emph{Statistica Sinica} 14.
\item Ray, K. and van der Vaart, A. W. (2020). Semiparametric Bayesian causal inference. \emph{Annals of Statistics} 48.
\item Reiers\o l, O. (1941). Confluence analysis by means of lag moments and other methods of confluence analysis. \emph{Econometrica} 9.
\item Reiers\o l, O. (1945). Confluence analysis by means of instrumental sets of variables. \emph{Arkiv f\"or Matematik, Astronomi och Fysik} 32A.
\item Richardson, T. S. and Robins, J. M. (2013). Single world intervention graphs: a primer. Working paper, University of Washington.
\item Ritov, Y., Bickel, P. J., Gamst, A. C. and Kleijn, B. J. K. (2014). The Bayesian analysis of complex, high dimensional models: can it be CODA? \emph{Statistical Science} 29.
\item Robins, J. M. (1986). A new approach to causal inference in mortality studies with a sustained exposure period. \emph{Mathematical Modelling} 7.
\item Robins, J. M. (1994). Correcting for non compliance in randomized trials using structural nested mean models. \emph{Communications in Statistics, Theory and Methods} 23.
\item Robins, J. M. and Ritov, Y. (1997). Toward a curse of dimensionality appropriate (CODA) asymptotic theory for semiparametric models. \emph{Statistics in Medicine} 16.
\item Robins, J. M., Hern\'an, M. A. and Brumback, B. (2000). Marginal structural models and causal inference in epidemiology. \emph{Epidemiology} 11.
\item Robins, J. M., Li, L., Mukherjee, R., Tchetgen Tchetgen, E. and van der Vaart, A. W. (2017). Minimax estimation of a functional on a structured high dimensional model. \emph{Annals of Statistics} 45.
\item Robins, J. M., Li, L., Tchetgen Tchetgen, E. and van der Vaart, A. W. (2008). Higher order influence functions and minimax estimation of nonlinear functionals. In \emph{Probability and Statistics: Essays in Honor of David A. Freedman}, Institute of Mathematical Statistics.
\item Robins, J. M., Rotnitzky, A. and Zhao, L. P. (1994). Estimation of regression coefficients when some regressors are not always observed. \emph{Journal of the American Statistical Association} 89.
\item Robins, J. M., Sued, M., Lei Gomez, Q. and Rotnitzky, A. (2007). Comment: performance of double robust estimators when ``inverse probability'' weights are highly variable. \emph{Statistical Science} 22.
\item Robins, J. M., Tchetgen Tchetgen, E., Li, L. and van der Vaart, A. W. (2009). Semiparametric minimax rates. \emph{Electronic Journal of Statistics} 3.
\item Robinson, P. M. (1988). Root N consistent semiparametric regression. \emph{Econometrica} 56.
\item Rockafellar, R. T. and Uryasev, S. (2000). Optimization of conditional value at risk. \emph{Journal of Risk} 2.
\item Ro\v{c}kov\'a, V. and Saha, E. (2019). On theory for BART. \emph{Proceedings of the 22nd International Conference on Artificial Intelligence and Statistics}, PMLR 89.
\item Ro\v{c}kov\'a, V. and van der Pas, S. (2020). Posterior concentration for Bayesian regression trees and forests. \emph{Annals of Statistics} 48.
\item Rosenbaum, P. R. (2002). \emph{Observational Studies}, second edition. Springer.
\item Rosenbaum, P. R. and Rubin, D. B. (1983). The central role of the propensity score in observational studies for causal effects. \emph{Biometrika} 70.
\item Rosenblatt, M. (1952). Remarks on a multivariate transformation. \emph{Annals of Mathematical Statistics} 23.
\item Rossi, P. E., Allenby, G. M. and McCulloch, R. (2005). \emph{Bayesian Statistics and Marketing}. Wiley.
\item Roth, A. E., S\"onmez, T. and \"Unver, M. U. (2004). Kidney exchange. \emph{Quarterly Journal of Economics} 119.
\item Roth, J., Sant'Anna, P. H. C., Bilinski, A. and Poe, J. (2023). What's trending in difference in differences? A synthesis of the recent econometrics literature. \emph{Journal of Econometrics} 235.
\item Rothstein, J. (2007). Does competition among public schools benefit students and taxpayers? Comment. \emph{American Economic Review} 97.
\item Rotnitzky, A., Smucler, E. and Robins, J. M. (2021). Characterization of parameters with a mixed bias property. \emph{Biometrika} 108.
\item Rubin, D. B. (1974). Estimating causal effects of treatments in randomized and nonrandomized studies. \emph{Journal of Educational Psychology} 66.
\item Rubin, D. B. (1978). Bayesian inference for causal effects: the role of randomization. \emph{Annals of Statistics} 6.
\item Rubin, D. B. (1981). The Bayesian bootstrap. \emph{Annals of Statistics} 9.
\item Rubin, D. B. (1984). Bayesianly justifiable and relevant frequency calculations for the applied statistician. \emph{Annals of Statistics} 12.
\item Rue, H., Martino, S. and Chopin, N. (2009). Approximate Bayesian inference for latent Gaussian models by using integrated nested Laplace approximations. \emph{Journal of the Royal Statistical Society, Series B} 71.
\item Sanov, I. N. (1957). On the probability of large deviations of random variables. \emph{Matematicheskii Sbornik} 42.
\item Sargan, J. D. (1958). The estimation of economic relationships using instrumental variables. \emph{Econometrica} 26.
\item S\"arndal, C. E., Swensson, B. and Wretman, J. (2003). \emph{Model Assisted Survey Sampling}. Springer.
\item Sarsons, H. (2015). Rainfall and conflict: a cautionary tale. \emph{Journal of Development Economics} 115.
\item Schwartz, L. (1965). On Bayes procedures. \emph{Zeitschrift f\"ur Wahrscheinlichkeitstheorie und verwandte Gebiete} 4.
\item Senn, S. (2023). Student and the Lanarkshire milk experiment. \emph{European Journal of Epidemiology} 38.
\item Shaked, M. and Shanthikumar, J. G. (2007). \emph{Stochastic Orders}. Springer, New York.
\item Shimizu, S., Hoyer, P. O., Hyv\"arinen, A. and Kerminen, A. (2006). A linear non Gaussian acyclic model for causal discovery. \emph{Journal of Machine Learning Research} 7.
\item Shpitser, I. and Pearl, J. (2006). Identification of joint interventional distributions in recursive semi Markovian causal models. \emph{Proceedings of the 21st National Conference on Artificial Intelligence}.
\item Sims, C. A. (2012). On an example of Larry Wasserman, round 2. Unpublished note, Princeton University.
\item Singh, R., Sahani, M. and Gretton, A. (2019). Kernel instrumental variable regression. \emph{Advances in Neural Information Processing Systems} 32.
\item Skilling, J. (2006). Nested sampling for general Bayesian computation. \emph{Bayesian Analysis} 1.
\item Smith, J. A. and Todd, P. E. (2005). Does matching overcome LaLonde's critique of nonexperimental estimators? \emph{Journal of Econometrics} 125.
\item Smith, J. Q. (2010). \emph{Bayesian Decision Analysis: Principles and Practice}. Cambridge University Press.
\item Smith, T. M. F. (1976). The foundations of survey sampling: a review. \emph{Journal of the Royal Statistical Society, Series A} 139.
\item Sparapani, R., Spanbauer, C. and McCulloch, R. (2021). Nonparametric machine learning and efficient computation with Bayesian additive regression trees: the BART R package. \emph{Journal of Statistical Software} 97.
\item Spirtes, P., Glymour, C. and Scheines, R. (2000). \emph{Causation, Prediction, and Search}, second edition. MIT Press.
\item Staiger, D. and Stock, J. H. (1997). Instrumental variables regression with weak instruments. \emph{Econometrica} 65.
\item stochtree developers (2025). \emph{stochtree}: stochastic tree ensembles for supervised learning and causal inference. Python package.
\item Stock, J. H. and Trebbi, F. (2003). Retrospectives: who invented instrumental variable regression? \emph{Journal of Economic Perspectives} 17.
\item Stolley, P. D. (1991). When genius errs: R. A. Fisher and the lung cancer controversy. \emph{American Journal of Epidemiology} 133.
\item Student (1908). The probable error of a mean. \emph{Biometrika} 6.
\item Student (1931). The Lanarkshire milk experiment. \emph{Biometrika} 23.
\item Student (1938). Comparison between balanced and random arrangements of field plots. \emph{Biometrika} 29.
\item Sugden, R. A. and Smith, T. M. F. (1984). Ignorable and informative designs in survey sampling inference. \emph{Biometrika} 71.
\item Sun, L. and Abraham, S. (2021). Estimating dynamic treatment effects in event studies with heterogeneous treatment effects. \emph{Journal of Econometrics} 225.
\item Suppes, P. (1970). \emph{A Probabilistic Theory of Causality}. North Holland.
\item Szekeres, G. (1958). Regular iteration of real and complex functions. \emph{Acta Mathematica} 100.
\item Talts, S., Betancourt, M., Simpson, D., Vehtari, A. and Gelman, A. (2018). Validating Bayesian inference algorithms with simulation based calibration. arXiv:1804.06788.
\item Theil, H. (1953). Repeated least squares applied to complete equation systems. Mimeograph, Central Planning Bureau, The Hague.
\item Thistlethwaite, D. L. and Campbell, D. T. (1960). Regression discontinuity analysis: an alternative to the ex post facto experiment. \emph{Journal of Educational Psychology} 51.
\item Thorgeirsson, T. E. et al. (2008). A variant associated with nicotine dependence, lung cancer and peripheral arterial disease. \emph{Nature} 452.
\item Thornton, R. L. (2008). The demand for, and impact of, learning HIV status. \emph{American Economic Review} 98.
\item Tian, J. and Pearl, J. (2000). Probabilities of causation: bounds and identification. \emph{Annals of Mathematics and Artificial Intelligence} 28.
\item Tibshirani, R. (1996). Regression shrinkage and selection via the lasso. \emph{Journal of the Royal Statistical Society, Series B} 58.
\item Titsias, M. (2009). Variational learning of inducing variables in sparse Gaussian processes. \emph{Proceedings of the 12th International Conference on Artificial Intelligence and Statistics}, PMLR 5.
\item Tobacco and Genetics Consortium (2010). Genome-wide meta-analyses identify multiple loci associated with smoking behavior. \emph{Nature Genetics} 42.
\item Trotter, H. F. and Tukey, J. W. (1956). Conditional Monte Carlo for normal samples. In \emph{Symposium on Monte Carlo Methods}, Wiley.
\item Truong, T. et al. (2010). Replication of lung cancer susceptibility loci at chromosomes 15q25, 5p15, and 6p21: a pooled analysis from the International Lung Cancer Consortium. \emph{Journal of the National Cancer Institute} 102.
\item Tsiatis, A. A. (2006). \emph{Semiparametric Theory and Missing Data}. Springer.
\item Tversky, A. and Kahneman, D. (1992). Advances in prospect theory: cumulative representation of uncertainty. \emph{Journal of Risk and Uncertainty} 5.
\item United States Department of Health, Education, and Welfare (1964). \emph{Smoking and Health: Report of the Advisory Committee to the Surgeon General of the Public Health Service}. Public Health Service Publication 1103.
\item van de Geer, S., B\"uhlmann, P., Ritov, Y. and Dezeure, R. (2014). On asymptotically optimal confidence regions and tests for high dimensional models. \emph{Annals of Statistics} 42.
\item van der Laan, M. J. and Rubin, D. B. (2006). Targeted maximum likelihood learning. \emph{International Journal of Biostatistics} 2.
\item van der Pas, S. L., Kleijn, B. J. K. and van der Vaart, A. W. (2014). The horseshoe estimator: posterior concentration around nearly black vectors. \emph{Electronic Journal of Statistics} 8.
\item van der Pas, S. L., Szab\'o, B. T. and van der Vaart, A. W. (2017a). Adaptive posterior contraction rates for the horseshoe. \emph{Electronic Journal of Statistics} 11.
\item van der Pas, S. L., Szab\'o, B. T. and van der Vaart, A. W. (2017b). Uncertainty quantification for the horseshoe, with discussion. \emph{Bayesian Analysis} 12.
\item van der Vaart, A. W. (1998). \emph{Asymptotic Statistics}. Cambridge University Press.
\item VanderWeele, T. J. (2015). \emph{Explanation in Causal Inference: Methods for Mediation and Interaction}. Oxford University Press.
\item VanderWeele, T. J. and Ding, P. (2017). Sensitivity analysis in observational research: introducing the E value. \emph{Annals of Internal Medicine} 167.
\item Vansteelandt, S. and Daniel, R. M. (2014). On regression adjustment for the propensity score. \emph{Statistics in Medicine} 33.
\item Vansteelandt, S. and Daniel, R. M. (2017). Interventional effects for mediation analysis with multiple mediators. \emph{Epidemiology} 28.
\item Vehtari, A., Gelman, A., Simpson, D., Carpenter, B. and B\"urkner, P.-C. (2021). Rank-normalization, folding, and localization: an improved $\widehat R$ for assessing convergence of MCMC. \emph{Bayesian Analysis} 16.
\item Viard, V. B., Gron, A. and Polson, N. G. (2015). Bayesian estimation of nonlinear equilibrium models with random coefficients. \emph{Applied Stochastic Models in Business and Industry} 31.
\item Viard, V. B., Polson, N. G. and Gron, A. (2001). Solving for market equilibrium using random coefficient random utility models. \emph{Computing in Economics and Finance 2001, No. 256}.
\item Villani, C. (2009). \emph{Optimal Transport: Old and New}. Springer.
\item von Neumann, J. and Morgenstern, O. (1944). \emph{Theory of Games and Economic Behavior}. Princeton University Press.
\item Wager, S. and Athey, S. (2018). Estimation and inference of heterogeneous treatment effects using random forests. \emph{Journal of the American Statistical Association} 113.
\item Wald, A. (1940). The fitting of straight lines if both variables are subject to error. \emph{Annals of Mathematical Statistics} 11.
\item Wang, S. S. (2000). A class of distortion operators for pricing financial and insurance risks. \emph{Journal of Risk and Insurance} 67.
\item Wang, Y., Broderick, P., Matakidou, A., Eisen, T. and Houlston, R. S. (2011). Chromosome 15q25 (CHRNA3-CHRNA5) variation impacts indirectly on lung cancer risk. \emph{PLoS ONE} 6.
\item Wasserman, L. (2004). \emph{All of Statistics}. Springer.
\item Weinberg, M. (2008). The price effects of horizontal mergers. \emph{Journal of Competition Law and Economics} 4.
\item Weitzman, M. L. (2009). On modeling and interpreting the economics of catastrophic climate change. \emph{Review of Economics and Statistics} 91.
\item Welch, B. L. and Peers, H. W. (1963). On formulae for confidence points based on integrals of weighted likelihoods. \emph{Journal of the Royal Statistical Society, Series B} 25.
\item Werden, G. J. (1996). A robust test for consumer welfare enhancing mergers among sellers of differentiated products. \emph{Journal of Industrial Economics} 44.
\item Werden, G. J. and Froeb, L. M. (1994). The effects of mergers in differentiated products industries: logit demand and merger policy. \emph{Journal of Law, Economics, and Organization} 10.
\item Woodward, J. (2003). \emph{Making Things Happen: A Theory of Causal Explanation}. Oxford University Press.
\item Wooldridge, J. M. (2016). Should instrumental variables be used as matching variables? \emph{Research in Economics} 70.
\item Wooldridge, J. M. (2020). \emph{Introductory Econometrics: A Modern Approach}, seventh edition. Cengage.
\item Wright, P. G. (1928). \emph{The Tariff on Animal and Vegetable Oils}. Macmillan.
\item Wright, S. (1921). Correlation and causation. \emph{Journal of Agricultural Research} 20.
\item Wright, S. (1934). The method of path coefficients. \emph{Annals of Mathematical Statistics} 5.
\item Wu, G. and Gonzalez, R. (1996). Curvature of the probability weighting function. \emph{Management Science} 42.
\item Wynder, E. L. and Graham, E. A. (1950). Tobacco smoking as a possible etiologic factor in bronchiogenic carcinoma. \emph{Journal of the American Medical Association} 143.
\item Yaari, M. E. (1987). The dual theory of choice under risk. \emph{Econometrica} 55.
\item Yakowitz, S., Krimmel, J. E. and Szidarovszky, F. (1978). Weighted Monte Carlo integration. \emph{SIAM Journal on Numerical Analysis} 15.
\item Yiu, A., Fong, E., Holmes, C. and Rousseau, J. (2025). Semiparametric posterior corrections. \emph{Journal of the Royal Statistical Society, Series B} 87.
\item Yoon, J., Jordon, J. and van der Schaar, M. (2018). GANITE: estimation of individualized treatment effects using generative adversarial nets. \emph{International Conference on Learning Representations}.
\item Yu, K. and Moyeed, R. A. (2001). Bayesian quantile regression. \emph{Statistics and Probability Letters} 54.
\item Zhai, P. S. and Ro\v{c}kov\'a, V. (2026). Horseshoe predictive inference. arXiv:2604.16661.
\item Zhang, C. H. and Zhang, S. S. (2014). Confidence intervals for low dimensional parameters in high dimensional linear models. \emph{Journal of the Royal Statistical Society, Series B} 76.
\item Zheng, X., Aragam, B., Ravikumar, P. and Xing, E. P. (2018). DAGs with NO TEARS: continuous optimization for structure learning. \emph{Advances in Neural Information Processing Systems} 31.
\item Zigler, C. M., Watts, K., Yeh, R. W., Wang, Y., Coull, B. A. and Dominici, F. (2013). Model feedback in Bayesian propensity score estimation. \emph{Biometrics} 69.
\item Zubizarreta, J. R. (2015). Stable weights that balance covariates for estimation with incomplete outcome data. \emph{Journal of the American Statistical Association} 110.
\end{list}
